% !TeX root = ../main.tex
% ─────────────────────────────────────────────────────────────
%  Abstract — English edition. 150–300 words answering the four
%  questions: problem / what we did / results / significance.
% ─────────────────────────────────────────────────────────────

\begin{abstract}
% (1) Problem
Many Vietnamese folk board games survive only by word of mouth among village
children, and they are slowly being forgotten. \textbf{Cờ Trận} is one such
game. It was played in Đồng Bùa hamlet, Tam Quan commune, Tam Đảo district,
Vĩnh Phúc province, using whatever could be picked up by the roadside as
pieces: stones, pebbles, twigs, leaves. To our knowledge it has never been
written down.
% (2) What we did
This paper records the complete rules of Cờ Trận for two, three and four
players, from the recollection of the author, who played it throughout
grades 5 to 9. We refer to it in English as \emph{Tactical Formation}. A
piece moves by jumping over an adjacent piece of any side, and the
jumped-over piece stays on the board. After landing, the moving piece
sandwiches the line of enemy pieces lying between it and a friendly piece,
and that whole line changes side. The board has two layers of positions,
vertices and centres: a vertex piece moves in 8 directions and a centre
piece in 4.
% (3) Results
We give a formal mathematical model of the two-player mode, stating a
\emph{28-piece conservation law} (capture $=$ change of side, never
removal). We also implement the game as open-source software, with an AI
based on Minimax search with Alpha-Beta pruning (the Paranoid variant for the
multiplayer modes). Experiments measure the branching factor and the think
time per move, and a two-player bot-vs-bot tournament shows that deeper
search wins consistently.
% (4) Significance
Through a written record, a formal model and software that anyone can play,
the paper helps keep a folk game from disappearing entirely.
\end{abstract}

\vspace{0.4em}
\noindent\textbf{Keywords:}
Vietnamese folk game;
traditional board game;
intangible cultural heritage;
abstract board game;
Paranoid Minimax;
multiplayer board game.
